\section{Two ranks on one marked arithmetic frame}

% BEGIN A68M1
\paragraph{A68M1. Returning the three outputs from the two determinant ideals.}
Fix one original frame $F$, its two setups $T_2,T_3$ and their
respective ideals $J_2,J_3$. Assume exactly
$T_i.\mathrm{ModIn}(J_i)$ for $i=2,3$. The coefficient rings of the
matrix spaces and the determinant rings of the two setups may differ.
By V68M4 choose a surjective comparison $\Phi_i$ for each, and
write $M_i=J_i S_i.B$. Each $\mathrm{ModIn.fin}$ supplies its own
weight algebra and finite quotient, so W67F2 gives
$\mathrm{FinLam}(M_i)$. The fields $\mathrm{n4}$ and $\mathrm{p1}$
give $\mathrm{P1}'$ in each rank by P68K7. In rank three use the
same $\mathrm{n4}$ and the field $\mathrm{p3}$ in that claim to
obtain $\mathrm{P23}'$. Finally $\mathrm{unr},\mathrm{dim}$ and
$\mathrm{seed}$ give the geometric seed on each of these very ideals
by S68D6. W67F4 now gives
\[
 T_2.\mathrm{GoodPt}(J_2),\qquad T_3.\mathrm{GoodPt}(J_3),
       \qquad \mathrm{FinPt}(O,F,3).
\]
This proves the original $\mathrm{AbsMain}$ implication. The two
good-point conclusions have determinant and character values before
passing to a fraction field. The finite-point assertion quantifies
over every original weight-zero test point.
% END A68M1

% BEGIN A68M2
\paragraph{A68M2. The dimension provided by a linked Hida weight algebra.}
Fix a linked frame over the CM field $L$, with positive-place set
$V$, and a chosen $v\in V$ as in the original arithmetic run.
Let $L^+$ be its maximal totally real subfield and $U=G_{L^+}$
in $G=G_{\mathbb Q}$. It is normal of index $[L^+:\mathbb Q]$
by abelianity. The frame chooses one representative $B_v$ of
each $U$-orbit of places above $\ell$ in the rational algebraic
closure. If $D_v$ is the stabilizer of $B_v$ in $G$, the local
degree supplied by the link is $[D_v:D_v\cap U]$.
The fibers of the projection $G/U\to U\backslash G/D_{v_0}$
have these sizes: on a double coset with representative $a$ the fiber
is the translate of the image of $aD_{v_0}a^{-1}$ by $aU$
in $G/U$, namely $\{adU:d\in D_{v_0}\}$. The places form the
conjugacy orbit of $B_{v_0}$. Thus the sum of local degrees is
\[
                          F.d=[L^+:\mathbb Q].
\]
Moreover the specified local group at $v$ contains an element
with nontrivial mod-$\ell$ cyclotomic character by the link's
$\mathrm{mu}$ field. It fixes $L^+$. Its place divides the modulus
$M$ because $\ell\mid M$, so the link's splitting criterion places
it in $\ker\eta$. The equality $G_L=U\cap\ker\eta$ then says it
fixes $L$. A primitive $\ell$th root in $L$ would force its
cyclotomic value to be one. This proves the required exclusion.

Given the original $\mathrm{LamEx}$ object, its weight ring in
rank $n$ is isomorphic to $\Lambda=O[[X_1,\dots,X_{nF.d}]]$.
Suppose $D/J$ is finite and torsion-free over $\Lambda$. For each
prime $Z$ minimal over $J$, the contraction of $Z/J$ to $\Lambda$
is zero. Indeed in the Noetherian ring $D/J$ a minimal prime is
associated, and every nonzero element of $\Lambda$ acts injectively.
The finite domain extension $\Lambda\hookrightarrow D/Z$ is integral.
Lying over, going up and incomparability give
$\dim(D/Z)=\dim\Lambda=1+nF.d$. This is the individual dimension
required for \emph{every} minimal prime of $J$.
% END A68M2

% BEGIN A68M3
\paragraph{A68M3. A local obstruction to a cyclotomic intertwiner.}
For a point $y$ of a residual determinant with the source
$\mathrm{V0ok}$ property, take $g$ in its distinguished local group
$D_0$ at which $\bar\varepsilon(g)\ne1$. All residual characters
agree on $g$, with common nonzero value $a$. Write $M_y$ for the
integral matrix of $y(g)$ over its curve DVR, $b=\varepsilon(g)$
and $p(X)=\det(X-M_y)$. Reduction of the point's determinant gives
$\bar p=(X-a)^n$. Thus $\bar M_y=aI+N$ for a nilpotent matrix
$N$ by Cayley--Hamilton. The matrix
$\bar b\bar M_y-aI=(\bar b-1)aI+\bar bN$ is invertible.
In particular $p(bM_y)$ has nonzero determinant after reduction
and is invertible over the fraction field. If $bM_yY=YM_y$ over
that field, polynomial intertwining gives
$p(bM_y)Y=Yp(M_y)=0$. Hence $Y=0$. This proves the precise
$\mathrm{N4}$ condition on every point of each supplied linked
arithmetic setup. It requires just the equality of the residual
eigenvalues at $g$, and permits nontrivial residual Jordan blocks.

Fix a realization $R=\mathrm{CReal}$ in rank two or three, at the
given comparison, frame, link, coefficient extension and twist.
Assume $\mathrm{ModUAlt},\mathrm{LamEx}$, the required coefficient
largeness, and the specialization implication
\[
 \forall J,\quad R.B_0.\mathrm{PointTheoremBaseObject}(\iota,J)
             \Longrightarrow R.T.\mathrm{IdealSpecializationPoint}(J).
\]
Use A68M2 and the specified Hida level equality to apply
$\mathrm{ModUAlt}$ to $R.B_0$ and the chosen weight object.
It supplies one ideal $J$. Its ring fields supply finiteness,
torsion-freeness and unramifiedness at the distinguished inertia,
so A68M2 supplies $\mathrm{dim}$ and its inertia field supplies
$\mathrm{unr}$ of $R.T.\mathrm{ModIn}$. The rank is positive
and $R.v0$ supplies precisely the residual hypothesis used above
for $\mathrm{n4}$. Take $\mathrm{p1}$ and $\mathrm{p3}$ from the
same modular ideal. Applying the displayed implication to its point
field gives $\mathrm{seed}$. This constructs all seven fields of
$R.T.\mathrm{ModIn}(J)$, and the classicality field
$\mathrm{arithmeticPointClassicalityDense}(\iota,J)$ is the
$\mathrm{end2}$ field of the very same ideal.
% END A68M3

% BEGIN A68M4
\paragraph{A68M4. The arithmetic point of one run.}
Keep $R$ and $J$, and assume its $\mathrm{GoodPt}$ and classicality
properties. Take any coefficient DVR $A$ and marked ordinary point
in the original $\mathrm{AutPtAlt}$ test, with Galois realization
$r$ irreducible. Take the compositum of $\operatorname{Frac}(A)$ and
$\operatorname{Frac}(c_1.O)$ inside $\overline{\mathbb Q}_\ell$,
where $c_1.O$ is the chosen coefficient DVR of $R$. Its ring of integers contains both
rings, since $A$ is finite over the original DVR. The inclusions
are local: their maximal ideals consist of elements of positive
valuation. This is an allowed coefficient-domain test. Extending
the normalized representation to this ring and extending the
frame by its prescribed residue map preserves its reduction,
ordinary characters and continuity. The irreducibility test for
$\mathrm{GoodPt}$ holds: a subspace stable under the matrix-frame
image would be stable under the given $r$. Consequently A68M1's
good-point output gives a homomorphism of the original determinant
ring into the enlarged DVR killing $J$. Composing with its inclusion
in $\overline{\mathbb Q}_\ell$ gives a point $x$. The equalities
of its determinant and characters are preserved by that inclusion.
After restriction by the Galois projection its determinant is the
full law of $r\otimes\kappa$, by the realization's twist identities.

Here $\kappa$ is the specified algebraic twist, with the original
$\mathrm{IsKa}$ hypothesis. At $v\mid\ell$ its local algebraic
weights, in the reciprocal convention of the source, are denoted by
$t_{v,\tau}$. The untwisted inertial character in position $i$ is
$\Psi^{w_i}$, with $w=(1,-1)$ or $(1,0,-1)$. It is
$\varepsilon^{e w_i}$ times a finite-order character: this follows
by writing $\Psi$ as cyclotomic to the power $e$ times its
Dirichlet factor. Global and local reciprocity for the algebraic
Hecke character of $\kappa$ give a finite-order factor times its
monomial in local embeddings on units. The cyclotomic character
has exponent $-1$ in this weight convention. Multiplication
therefore makes $x$ locally algebraic of weights
\[
                         t_{v,\tau}-e w_i .
\]
These strictly increase with $i$ since $e>0$. The open-kernel
factor is allowed in the original local weight predicate. The point
$x$ is thus arithmetic, represents $r\otimes\kappa$, and kills
$J$. Scalar twisting preserves irreducibility. Applying the
classicality field of $J$ to this particular point gives
$\mathrm{AutOffN}(R.s.v_0,\iota,r\otimes\kappa)$ with its original
window condition. Only the distinguished place of this run is excluded.
% END A68M4

% BEGIN A68M5
\paragraph{A68M5. Two runs remove both excluded places.}
Assume the original two-run hypotheses, including the fixed ample
coefficient field, the two prescribed twists, the modularity and
weight-algebra inputs of A68M3 and global reciprocity. In each rank
$i=2,3$ the hypothesis supplies a realization whose auxiliary
rational prime avoids any specified nonzero integer. For every such
realization assume the specialization implication of A68M3. Choose
first the realizations $R_2,R_3$ avoiding $1$. In each rank choose
a second realization avoiding the first auxiliary prime in that
rank. The two excluded places in each rank are distinct: the
distinguished-place norm in each realization is its auxiliary prime.

Use A68M3 on each realization. On each pair of ranks A68M1 now
gives its good points and $\mathrm{FinPt}$ on the same extended
frame. Keep the first finite-point conclusion. In either rank fix
an arbitrary original arithmetic test $r$ as in A68M4. The two
realizations give $\mathrm{AutOffN}$ for exactly the same
$r\otimes\kappa$ and $\iota$ with distinct exclusions.
HIA5 gives $\mathrm{AutCsd0}$ and the conditional
$\mathrm{WitA}$ for this representation, including one center
before all eligible places for each fixed attached cusp.
This is the original $\mathrm{IsAutCsdA}$. As the test was
arbitrary, we get both $\mathrm{AutPtAlt}$ assertions and
$\mathrm{FinPt}$ simultaneously, as in the source's
$\mathrm{twoRuns\_alt}$. RTC24 descends the finite assertion
when used in M65A4. To apply this conditional result to the
actual marked problem one must still produce the modularity and
specialization data, the indicated realizations and twists on the
given frame, and the polarized descent required by M65A4.
% END A68M5
